\chapter*{Theoretical Background: Cantilever Beam}
\addcontentsline{toc}{chapter}{Theoretical Background}
The fundamental objective of modal analysis is to derive the natural frequencies and mode shapes of a structure. Using a cantilever beam as a baseline analytical model, the standing wave solution is derived by applying appropriate boundary conditions (null displacement and rotation at the clamped end; moment and vertical force equilibrium at the free end). This yields an eigenvalue problem from which the structure's exact natural frequencies and continuous shape functions are computed.

Experimentally, Frequency Response Functions (FRFs) are measured for different input-output position pairs. By employing a multi-mode curve fitting method, these FRFs are approximated numerically to minimize the least-square error with the real experimental data. The modal parameters, such as natural frequencies and modal damping ratios, are extracted directly from the fitted FRFs. Finally, the mode shapes are reconstructed by evaluating the imaginary part of the FRFs at the resonant frequencies. This core theoretical foundation is directly applicable to the modal identification of much more complex physical systems.
